\section{Experimental Setup}\label{sec:experiments}

\subsection{Development subset and full-scale protocol}

The development corpus contains 150 clips and all 128 body configurations,
for 19,200 variants. These clips are reserved within the full training split
and are not a held-out generalization benchmark. Small-scale experiments
investigate architecture and reference-processing choices with a fixed
simulation budget. The full-scale run uses the 18,855 training clips in
Table~\ref{tab:dataset_counts}, with validation membership fixed throughout.

The selected full-scale experiment is slot/type attention trained directly on
refined references. Its code retains a historical \texttt{stage2} name, but
this run does not use neutral pretraining followed by shape transfer.
Table~\ref{tab:training_parameters} records the shared optimization settings
and the actual environment/batch sizes used after the initial memory-related
launch adjustments. Run-specific resolved configurations take precedence
over older launch examples.

\begin{table}[t]
\centering
\small
\caption{Training configuration and budgets. Environment and minibatch counts
are per GPU. The full-scale frame count is a recorded stopping snapshot,
not a budget matched to the development runs.}\label{tab:training_parameters}
\begin{tabularx}{\linewidth}{@{}lYY@{}}
\toprule
Setting & Corrected 150-clip runs & Full-scale selected run \\
\midrule
Parallel GPUs per policy & 1 & 6 \\
Environments per GPU & 4,096 & 2,048 \\
PPO minibatch size per GPU & 16,384 & 8,192 \\
Rollout length & 32 control steps & 32 control steps \\
Mini-epochs per update & 1 & 1 \\
Actor / critic learning rates &
$2\times10^{-5}$ / $10^{-4}$ & $2\times10^{-5}$ / $10^{-4}$ \\
Discount / GAE / ratio clip & 0.99 / 0.95 / 0.2 & 0.99 / 0.95 / 0.2 \\
Gradient norm clip & 50 & 50 \\
Evaluation cadence & 200 epochs & 256 epochs \\
Training seed & 0 for current five runs & Recorded in saved run configuration \\
Reference residency & All 150 clips & 256, later 128 clips per GPU \\
Training frames & 786,432,000 planned per run & Approximately $11.07\times10^9$ recorded \\
Status of reported evidence & Comparisons pending & Monitoring at epoch 28,159 \\
\bottomrule
\end{tabularx}
\end{table}

\subsection{Architecture and refinement controls}\label{sec:ablation_design}

The corrected comparison consists of the five cases in
Table~\ref{tab:ablations}. Cases A--D use the same refined corpus and temporal
input categories. They test flat temporal aggregation, attention, slot/type
embeddings, and actor-only morphology modulation. Case E uses the same
architecture as C with unrefined references.

\begin{table}[t]
\centering
\small
\caption{Corrected 150-clip comparison. One seed per case is currently running;
this is an experiment design, not a ranking of completed results.}\label{tab:ablations}
\begin{tabularx}{\linewidth}{@{}llYY@{}}
\toprule
Case & References & Architecture & Comparison purpose \\
\midrule
A & Refined & Wide temporal MLP; beta concatenation &
Temporal MLP baseline. \\
B & Refined & Basic temporal attention &
Attention versus flattened temporal input. \\
C & Refined & Attention with slot/type embeddings &
Explicit temporal/source identity. \\
D & Refined & C + actor-only adaLN-Zero &
Morphology modulation with critic unchanged. \\
E & Unrefined & Same architecture as C &
Reference refinement at fixed architecture. \\
\bottomrule
\end{tabularx}
\end{table}

All five current runs use seed 0, one A40 GPU per policy, the same frame
budget, and the same evaluation cadence and shape-sampling seed. Morphology
concatenation into both heads is retained in D; adaLN-Zero is additional
actor conditioning, not a replacement for the critic or a two-stage schedule.

The models are not asserted to have equal parameter counts. Counts, wall time,
and GPU-hours must accompany the completed comparison, and repeated training
seeds are needed before making claims about reliability or statistical
significance. For C versus E, final controller evaluation must use identical
clip/shape pairs and a common reference version. Their training-time metrics
use different targets and are not by themselves a controlled measure of
refinement benefit. Cross-evaluation against both reference versions can
separate target fidelity from ease of tracking.

Earlier small-scale runs influenced the choice to begin full-scale slot/type
training. The corrected reruns were launched later to reassess that choice;
we do not present them as completed model selection preceding the full run.

\subsection{Morphology-matched evaluation and checkpoint scope}\label{sec:evaluation_protocol}

Each reference is rolled out only in an environment using its exact asset
identity. A selected reference therefore cannot be assigned to a different
body simply because an environment slot is available. Multiple rollout
batches may be needed when several references require the same body.

Routine monitoring selects one variant per clip. With seed $s=42$, stable
clip identity $c$, panel index $p$, and $N_c$ available variants, the choice is
\begin{equation}
    k(c,p) =
    \left(
       \operatorname{uint64}_{\mathrm{big}}
       \bigl(\operatorname{SHA256}(s\!:\!c)_{0:8}\bigr)
       +p
    \right)\bmod N_c .
    \label{eq:shape_panel}
\end{equation}
Here $s\!:\!c$ is the UTF-8 string formed as \texttt{seed:clip\_id}.
For scheduled monitoring, $p$ is the integer quotient of the checkpointed
epoch and evaluation cadence. This produces deterministic staggered shape
choices that rotate across evaluations and remain reproducible after resume.
The validation \emph{clip list} is fixed; the selected \emph{shape panel}
changes. Thus successive scores are not repeated measurements on an
identical set of clip/shape pairs.

Evaluation uses the mean action, reference initialization, and no gradient
tracking. It disables training termination and records errors over the
valid reference horizon, capped at 600 control steps (20\,s at 30\,Hz).
Success therefore refers to the evaluated horizon, not necessarily the
entire duration of a longer source clip.

Final model comparisons should freeze the checkpoint, clip/shape panel,
reference version, simulator configuration, and seed. Validation is available
for model selection and has been repeatedly monitored; it is not an untouched
test set. The independent test manifest remains outside training, and no
test result is reported in this draft. The downloaded epoch-28,170
\texttt{last.ckpt} is also distinct from the epoch-28,159 monitoring snapshot
in Table~\ref{tab:full_results}; a new evaluation of that checkpoint must be
reported separately.

\subsection{Metrics and aggregation}\label{sec:metrics}

For each evaluated clip/shape pair $i$, define its per-frame mean body
position error
\begin{equation}
    d_{i,t} = \frac{1}{B}\sum_{j=1}^{B}
                 \|p_{i,t,j}-\hat{p}_{i,t,j}\|_2 .
\end{equation}
The configured success indicator is
\begin{equation}
    S_i = \mathbf{1}\!\left[
             \max_{t<T_i}d_{i,t}\leq 0.5\ \mathrm{m}
          \right],
    \label{eq:success}
\end{equation}
where $T_i$ is the evaluated horizon. This threshold is deliberately distinct
from the smoother training reward. Passing it does not imply fine-grained
pose accuracy or biomechanical realism.

We also report time-averaged mean body-position error, mean global
body-rotation geodesic error, and the time average of the largest body-position
error within each frame. The last quantity is the logged
\emph{mean maximum-joint error}; it is not the maximum over the entire dataset.
Global body-rotation error is not a local joint-angle/DOF error.

Within an evaluator rank, errors are averaged over valid frames per motion
and then over motions, including failed rollouts. The distributed training
logger aggregates scalar summaries across ranks; the monitoring table
retains those logged values rather than reconstructing exact pooled counts.
Final offline comparisons should aggregate saved per-motion records over
unique clip/shape pairs, avoiding unequal weighting of partially filled
batches or duplicate padded rollouts.

Trajectory smoothness uses the implementation's windowed normalized-jerk
score, with the definition and numerical regularizers in
Appendix~\ref{app:smoothness}. Lower scores indicate lower jerk under that
fixed convention. The accompanying high-jerk statistic is the percentage
of windows in which any body exceeds threshold 6,500; the logger calls it
a frame percentage, but it is window-based.

Action increments require care. The current evaluator stores raw policy
actions, before Equation~\eqref{eq:action_mapping}. Their differences are
policy-output changes, not physical joint-angle changes, even where logger
field names label them radians or degrees. We omit those mislabeled angular
values from the physical interpretation and from Table~\ref{tab:full_results}.
PD-target changes, applied torques, and measured joint motion must be
computed separately.

\subsection{Generalization axes and compute accounting}

Held-out clips on trained bodies, seen clips on unseen bodies, and held-out
clips on unseen bodies are different evaluations. The current validation
snapshot addresses only the first. A single rotating panel does not
measure all 128 variants of every clip. Future cross-shape comparisons must
report the actual body identities and number of evaluated pairs.

We retain W\&B logs, resolved configurations, checkpoint epochs, and local
results for compute accounting. Small-scale development and full-scale
training costs should be reported separately, including resumed segments and
failed launches when discussing practical cost. GPU-hours, rather than
epoch counts alone, quantify this budget.
\pending{Consolidate per-run parameter counts and GPU-hours from the stored
logs, and add repeated-seed comparisons. No superior compute-efficiency claim
is made from the current evidence.}
